\KOMAoption{paper}{a4, portrait}
\areaset{267mm}{180mm}
\fancyheadoffset{0pt}

\newgeometry{
  top=14mm,
  bottom=11mm,
  right=10mm,
  left=22mm,
}

\begin{landscape}
\setlength{\footskip}{15.5pt}
\color{unimain}{\section[(обязательное) Схема электрическая принципиальная]{(обязательное)\\Схема электрическая принципиальная}\label{app:a3}}
\color{black}\vskip 3mm

\newcount\figcounter
\figcounter=1
\loop
  \IfFileExists{\apppath/Приложение. Схема Э3/_latex/img/СЭП_\the\figcounter.pdf}{%
    \begin{figure}[H]
      \centering
      % Безопасные пределы: 95% от доступной ширины и высоты
      % Это гарантирует, что рисунок НИКОГДА не вылезет за поля
      \ifnum\figcounter=1
        \def\maxdim{0.8}% Для первой схемы чуть больше запаса
      \else
        \def\maxdim{0.95}% Для остальных — максимум
      \fi
      
      \includegraphics[
        width=\maxdim\linewidth,   % \linewidth в landscape = ширина альбомного листа
        height=\maxdim\textheight, % \textheight в landscape = высота альбомного листа
        keepaspectratio,
        angle=0
      ]{%
        \apppath/Приложение. Схема Э3/_latex/img/СЭП_\the\figcounter.pdf
      }%
      
      \captionof{figure}{Схема электрическая принципиальная (часть \the\figcounter)}%
      \label{appvo:fig_\the\figcounter}%
    \end{figure}%
    \advance\figcounter by 1
  }{%
    \figcounter=0 % файл не найден — завершаем цикл
  }%
\ifnum\figcounter>0 \repeat

\clearpage
\end{landscape}
\restoregeometry